\include{zLaTex-Preamble}

\title{Project 4 Design Writeup} 
\author{Joshua Miller}
% \date{}

\begin{document}
\rhead{\fancyplain{}{Joshua Miller}}
\maketitle

\renewcommand{\t}[1]{{\tt\sc #1}}

\section{General Design}

{\bf Modules}
\begin{itemize}
\item util.h -- contains macros and definitions for error handling and debugging
\item hash.h -- contains definitions for each of the four hash table types 
\item tests.c -- contains the tests for the entire package
\item app.h -- contains the different hash performance tests
  \t{Serial, Parallel, ParallelNoLoad}.
\item experiment\{1,2,3\}.c -- contains the code to run the experiments.
\end{itemize}

\section{General}

During the last project, I ran into debugging problems due to a bit of
carelessness in my polymorphic implementation of the lock structures.
Even with my bad experience here, I will try and proceed with the hash
tables by wrapping them into a polymorphic data structure.  Each will
have functions of the form:

\lang{C}
\begin{code}

/* Thread pool definitions */
typedef void *(*loop_t)(void*);

typedef struct thread_pool_t thread_pool_t;

typedef struct thread_t {
    pthread_t self;
    thread_pool_t *pool;
} thread_t;

typedef struct thread_pool_t{
    thread_t *threads;
    int size;
    int start;
    int stop;
} thread_pool_t;

/* The top level hash table abstraction */
typedef struct hasht_t hasht_t;   

/* Constructor for new hash table */
hash_t *hasht_new(hasht_type_t type, int capacity, int expected_threads);

/* Pointer types for each table interaction member */
typedef int (*hasht_free_f)    (hasht_t *);
typedef int (*hasht_get_type_f)(hasht_t *);
typedef int (*hasht_resize_f)  (hasht_t *);
typedef int (*hasht_add_f)     (hasht_t *, void *);
typedef int (*hasht_remove_f)  (hasht_t *, void *);
typedef int (*hasht_contains_f)(hasht_t *, void *);

/* Enumeration to be contained in hasht object */
typedef enum hasht_type_t 
{
  LOCKING,
  LOCK_FREE,
  LINEAR_PROBE,
  AWESOME
} hasht_type_t;

struct hasht_t 
{
  /* Type specifier */
  hasht_type_t type;

  long count;
  long capacity;

  /* Member Function Pointers */
  hasht_add_f        add;
  hasht_contains_f   contains;
  hasht_remove_f     remove;
  hasht_resize_f     resize;

  /* Member types */
  hasht_locking_t      locking;
  hasht_lockfreec_t    lock_free;
  hasht_linearprobe_t  probing;
  hasht_awesome_t      awesome;
}; 

\end{code}

The constructor for the hash table is responsible for assigning the
type to the object, as well as its respective functions. 


\section{Lock-based Closed-address}

The table contains the following
\begin{itemize}
\item an array of linked lists of size \t{Capacity}
\item an array of integers specifying the length of each bucket
\item an array of read locks of size \t{nlocks}
\item an array of write locks \t{nlocks}
\item an integer \t{Capacity}
\item an integer \t{Count} specifying how many packets are in the table
\end{itemize}

Because we can stripe the hash table, the locks are held in parallel
array, such that the lock for item $i$ can be accessed by calculating
$i \mod$\t{nlocks}.

Editing (writing, resizing) the hash table will require acquisition of
both the respective write and read locks. Reading (contains) from the
hash table will require acquisition of only the read locks.

\section{Lock-Free-Contains Closed-Address}

The table contains the following
\begin{itemize}
\item an array of lock free queues lists of size \t{Capacity}
\item an array of integers specifying the length of each bucket
\item an global lock to lock the state on resize
\item an array of write locks \t{nlocks}
\item an integer \t{Capacity}
\item an integer \t{Count} specifying how many packets are in the table
\end{itemize}

Each call to \t{add()} will increment the total count, and the bucket
count.  It will also check to make sure that the bucket length is not
greater than the threshold.  If it is, it will resize the hash
table. The main idea is that we must create a linearization point at
the insertion of a new element in the bucket that we can make atomic.

\section{Linearly Probed Open-Address}

The table contains the following
\begin{itemize}
\item a single array for elements
\item a last element in array index
\item a size
\end{itemize}

Each call to remove or add will step down the list with an interval of
1.  This means that it will check every element in the list for the
correct tag. If we have reached the limit set as the number of
occupied spaces, we atomicly insert the item. If there is no room for
the data, we will acquire the global lock and resize the list.


\section{My Design}

Because any custom design will be based off of the lock-free or
linear-probe, I will postpone the design of this until after
performance testing of both of the locks and make an attempt to modify
whichever demonstrates better performance.

\section{Tests}

Because of the polymorphic nature of this design, the correctness
testing will be easily accomplished through composite tests over the
functions
\{add(), remove(), contains(), resize()\} that can be applied to each
hash table and checking a predetermined state.  

The performance tests are nearly given. There is boilerplate for the
parallel test and the initialization of the hash tables is outlined
above. The \t{ParallelNoLoad} will be a reimplementation of the
parallel packet test with no load.


\section{Experiments}

\subsection{Dispatcher Rate}

This experiment involves running the parallel test with no load and
minimal packet size for the number of threads equal to the number of
cores. This will present us with the rate at which the dispatcher
thread can enqueue the packets onto the lamport queues. I predict that
the dispatcher rate will be based upon the relative performanc of
the \t{hashgenerator}, which I believe will be slightly slower, and
therefore the performance will fall just below the performance of the
dispatcher rate in previous projects.


\subsection{Parallel Overhead}

Configuration: $n = 1$, $W = 4000$ 

This experiment involves running the parallel test with each of the
parallel hash tables, but on a single thread.  In comparison with the
performance of the serial hash table the speedup will be less than
1. I predict that for all hash tables, the speedup will be lower for
the heavy write runs.  This is because the writing has an added
complexity relative to the serial version to handle concurrency that
the reading does not.

\subsection{Speedup}

Configuration: $n \in \{1, 2, 4, 8, 16, 32\}$, $W = 4000$

This experiment involves running a scaling test on all the parallel
hash packet implementations.  I predict that our program will be able
to acheive 80\% of ideal speedup, or just over 12 times the serial
version on 16 cores.  









\end{document}
